\documentclass[12pt]{book}
\usepackage{fontspec}
\setmainfont{Noto Serif CJK JP}
\usepackage[a4paper,textwidth=13cm,top=2.5cm,bottom=2.5cm]{geometry}
\usepackage{titlesec}
\emergencystretch=3em
\tolerance=2000

\title{One Word, Not Four\\\large Japanese Grammar in Images}
\author{Markus Bauernfeind}
\date{2026}

\begin{document}
\pagestyle{plain}
\maketitle
\tableofcontents

\chapter{Preface}

Most textbooks explain Japanese grammar by forcing it into English categories: particles get explained as prepositions, verb forms as tenses, constructions as translation patterns. This works for a while — until it doesn't, because Japanese is built on a logic of its own that doesn't map cleanly onto English.

This book tries something different: to show the image a native Japanese speaker holds in mind when choosing a particle or a form. に isn't simply "to," "at," or "in" — it's an anchor point something attaches to. を isn't just "the object marker" — it's a stretch of ground being fully covered. Once you understand this inner image, you stop needing to memorize which particle "goes with this situation" — you recognize why it's the natural choice.

\section*{What This Book Is Not}

It isn't a linguistic treatise aiming at completeness. Language is complicated, and a learner's image doesn't need to account for every academic edge case. The images presented here explain the most common and most instructive cases of a given particle or form — not necessarily every last peripheral use.

Nor is it a collection of final, "proven" truths. Even professional linguists don't fully agree on how to capture the core meaning of certain particles — に and は among them. The images in this book are carefully checked, working simplifications for learning — not claimed, settled linguistic facts. Where a chapter runs into a genuine limit or open question, that's stated openly rather than smoothed over.

\section*{How the Chapters Are Built}

Each chapter is devoted to one particle or form and centers on the image: the mental picture that comes naturally to a speaker choosing that form. Examples illustrate the image with concrete sentences. Where there's a real, instructive edge case, it's included — not to unsettle, but because such cases often show most clearly how the image actually works. Where a word's origin helps understanding, it appears as a short note at the end of the chapter.

Not every chapter needs every one of these building blocks. Some particles have an interesting edge case, others don't; some have a revealing etymology, others don't. No chapter is padded out artificially just to follow a template.

\section*{Who This Book Is For}

For people learning Japanese — not for a linguistics audience. Whether as a companion to a course, a reference, or for self-study, the goal stays the same: to build a genuine feel for why Japanese works the way it does — not to memorize by rote, but to understand.

\section*{A Note on How This Book Was Made}

Claude (Anthropic) was used as an AI assistant in preparing this book — for research, language drafting, and counterexample checking of the images presented. All concepts, decisions, and the final judgment on content are the author's.


\chapter{に — The Anchor Particle}

\section*{The Image}

Picture に as a finger pointing at a \textbf{fixed point} — a place, a moment in time, a goal, a state. に doesn't mark \emph{movement through} something; it marks \emph{the point something is aimed at or anchored to}.

English speakers often reach for "to," "at," "in," or "on" depending on the situation — four different words for what Japanese treats as a single, unified concept: something gets anchored to a specific point.

\section*{The Main Cases}

\textbf{1. Location of existence (not of action)}
\par\quad 東京に住んでいます。\\
\par\quad I live in Tokyo.\\

The point at which something \emph{is}, not the point it moves through. Compare で (the location of an action) — that's the stage; this is the anchor.

\textbf{2. Point in time}
\par\quad 7時に起きます。\\
\par\quad I get up at 7 o'clock.\\

Time is treated exactly like place: a fixed point something is anchored to.

\textbf{3. Goal/direction of movement}
\par\quad 学校に行きます。\\
\par\quad I'm going to school.\\

The destination of the movement — not the path there, but the point the movement is aimed at.

\textbf{4. Recipient/direction of an action}
\par\quad 友達に手紙を書きます。\\
\par\quad I'm writing a letter to my friend.\\

The action is aimed at a point — the same mechanism as with place and time, just with a person as the anchor.

\textbf{5. Change of state (the goal of a transformation)}
\par\quad 先生になりたいです。\\
\par\quad I want to become a teacher.\\

The state one "arrives" at — a change of state has a goal point, just like a movement does.

\section*{An Interesting Edge Case}

Because に is itself \textbf{direction-neutral} (the actual direction comes from the verb, not from に itself), the same sentence can sometimes carry two readings.

The sentence 議長に推薦された can mean "I was recommended \emph{to become} chairman" (に as goal) — or "I was recommended \emph{by} the chairman" (に as the source of the action). Both are grammatically valid; context has to decide. This isn't a flaw in the image — it's a direct consequence of it: if に only marks the point of reference and not the direction itself, genuine ambiguity is bound to appear whenever context doesn't supply a clear direction.

\bigskip\hrule\bigskip

\emph{For the curious:} に traces back to an Old Japanese element that originally marked a static location — in contrast to へ, which historically expressed motion toward. This root in the idea of an "anchor point" explains why such different-looking uses could grow out of に.


\chapter{で — The Boundary Particle}

\section*{The Image}

Picture で as a drawn boundary — a circle around an area, a line marking a limit, a point where one state tips into another. で says: \emph{here is the boundary within which, up to which, or at which something happens.}

Unlike に (a point something is anchored to), で always marks a \textbf{boundary} — sometimes as an area (a space within which something takes place), sometimes as a limit (an upper bound), sometimes as a tipping point (the switch from one state to the next).

\section*{The Main Cases}

\textbf{1. Location of an action}
\par\quad 公園で遊びます。\\
\par\quad I play in the park.\\

The park is the bounded area within which the action takes place — not simply where something \emph{exists} (that would be に), but the stage on which something \emph{happens}.

\textbf{2. Means/instrument}
\par\quad ペンで書きます。\\
\par\quad I write with a pen.\\

The pen stakes out the bounds within which writing becomes possible — the means as a bounded space that enables the action.

\textbf{3. Material}
\par\quad 紙で作ります。\\
\par\quad I make it out of paper.\\

The starting boundary something emerges from.

\textbf{4. Total amount or deadline}
\par\quad 1000円で買いました。\\
\par\quad I bought it for 1000 yen.\\

\par\quad 一週間でできます。\\
\par\quad It'll be done within a week.\\

An upper bound — for money, a limit; for time, a deadline.

\textbf{5. Cause}
\par\quad 地震で電車が止まりました。\\
\par\quad Because of the earthquake, the trains stopped.\\

Here the boundary is a \textbf{tipping point}: the moment one state (trains running) flips into another (trains stopped). The same principle carries あとで ("afterward," the tipping point between what's finished and what's beginning) and causal ので.

\textbf{6. Scope of comparison}
\par\quad クラスの中で一番背が高いです。\\
\par\quad He's the tallest in the class.\\

中で ("within the middle") explicitly bounds the group within which the comparison holds — the same boundary-setting function as the other cases, just especially visible here.

\bigskip\hrule\bigskip

\emph{For the curious:} で goes back to にて, an older combination of に and the particle て. Today's range of uses — location, means, material, limit, cause — traces back to a single idea: a set boundary that frames an action.


\chapter{へ — The Direction Particle}

\section*{The Image}

Picture へ as an arrow still \textbf{on its way} — not the point of arrival, but the movement itself, the process leading there. While に marks the point of arrival (the result), へ marks the process, the path, the moving-toward.

These two particles often overlap in everyday speech — 東京に行く and 東京へ行く both mean "I'm going to Tokyo." But the subtle difference remains: に emphasizes \emph{where} one arrives, へ emphasizes \emph{that} one is moving there.

\section*{The Main Cases}

\textbf{1. Direction of movement}
\par\quad 東京へ行きます。\\
\par\quad I'm heading toward Tokyo.\\

The focus is on the movement itself, less on the exact destination.

\textbf{2. Abstract development}
\par\quad 夏へと向かっています。\\
\par\quad We're heading toward summer.\\

Even without physical movement, the image holds: a process, a development moving toward something — "one step further" (一歩先へ) works the same way.

\textbf{3. Recipient of an action}
\par\quad 恋人へ手紙を書きます。\\
\par\quad I'm writing a letter to my sweetheart.\\

Here へ doesn't just mark the recipient — it emphasizes the \emph{process}: the letter is on its way, the act of writing is directed toward the person, not just a single point of contact.

\section*{An Interesting Edge Case}

With direct, immediate physical contact, へ feels unnatural:

\par\quad 猫が私に噛みつきました。\\
\par\quad The cat bit me.\\

Only に works here, not へ. This fits the image exactly: a bite isn't a process unfolding over a distance — it's a discrete point of contact, immediate. Where there's no unfolding, no moving-toward, へ has nothing to mark.

\bigskip\hrule\bigskip

\emph{For the curious:} へ goes back to the noun 辺 ("area, vicinity" — as in 道の辺, "roadside"). Out of this spatial, area-based meaning the direction particle developed over time — and later increasingly encroached on に's territory (the point of arrival), which explains the overlap between the two particles today.


\chapter{を — The Traversal Particle}

\section*{The Image}

Picture を as a stretch of ground being fully covered, worked through, or left behind. を doesn't simply mark "the object" as a grammatical category — it marks whatever the action \textbf{fully passes through, fully takes hold of, or fully leaves behind.}

Something marked with を gets thoroughly affected or transformed by the action. A place marked with を gets fully traversed by the movement. A starting point marked with を gets fully departed from by the action. Always the same idea: the full extent of something is covered.

\section*{The Main Cases}

\textbf{1. Object of an action}
\par\quad 小説を書きます。\\
\par\quad I'm writing a novel.\\

The novel comes into being entirely through the action — captured by を as the fully worked-through domain of the action.

\textbf{2. Traversed place}
\par\quad 山道を行きます。\\
\par\quad I'm walking along the mountain path.\\

Here を is the path itself, fully covered by the movement — not the destination (that would be に or へ), but the stretch itself.

\textbf{3. Point of departure}
\par\quad 大学を卒業します。\\
\par\quad I'm graduating from university.\\

University gets left behind as a fully completed phase — the same traversal idea, just at the end rather than the middle of the movement.

\textbf{4. Extent of time}
\par\quad 長い年月を過ごします。\\
\par\quad I'm spending long years (doing this).\\

Time, too, can be thought of as a stretch that gets covered.

\section*{A Sharp Contrast with で}

海を泳ぐ and 海で泳ぐ look almost identical but mean different things:

\par\quad 海を泳ぎます。\\
\par\quad I'm swimming across the sea.\\

\par\quad 海で泳ぎます。\\
\par\quad I'm swimming in the sea.\\

With を, the entire sea is thought of as a stretch to be covered — fully traversed. With で (the boundary particle), the sea is just the area \emph{within which} one swims locally, with no claim to covering it fully. Two particles, two completely different images — each staying true to its own picture.

\section*{A Second Contrast: を and に with Coercion and Permission}

In certain causative sentences (making someone do something), there's a choice between を and に — and the difference is exactly what you'd expect:

\par\quad 子供を勉強させます。 → with を: more like coercion — the child is "pushed all the way through" studying.\\
\par\quad 子供に勉強させます。 → with に: more like permission — the child is allowed to study.\\

を pushes all the way through; に merely permits. Two independent particle images that, together, explain a real difference in meaning.


\chapter{が — The Identification Particle}

\section*{The Image}

Picture が as a finger pointing at something for the first time: "This one!" が identifies an entity not yet established as the topic for the listener, and makes it the direct bearer of whatever the sentence goes on to say.

\section*{The Main Cases}

\textbf{1. Subject of a newly noticed fact}
\par\quad 雨が降っています。\\
\par\quad It's raining.\\

The rain is being registered and named as a fact right now — no prior mention needed.

\textbf{2. Object with verbs of state, feeling, and perception}
\par\quad サッカーが好きです。\\
\par\quad I like soccer.\\

\par\quad 音楽が聞こえます。\\
\par\quad I can hear music.\\

Here too, が identifies the bearer of the state or perception — grammatically an "object," but following the same identification logic.

\textbf{3. Existence}
\par\quad 猫がいます。\\
\par\quad There's a cat.\\

が identifies what exists.

\textbf{4. Exclusive identification}
\par\quad あいつが犯人です。\\
\par\quad \emph{He's} the culprit (and no one else).\\

Here が sharpens: not just "a new fact," but "this one and no other" — a more emphatic variant of the same basic function.

\textbf{5. Question words}
\par\quad 誰が来ますか。\\
\par\quad Who's coming?\\

Question words almost always take が — which makes sense, since a question word is by definition unknown and can't already be an established starting point.


\chapter{は — The Singling-Out Particle}

\section*{The Image}

Picture は as a hand lifting one element out of the sentence and placing it center stage: "This is what we're talking about now." は isn't an ordinary case particle like が or を — it often replaces them. Instead, は \textbf{singles out} an element and makes it the starting point for everything that follows.

This singling-out has two faces, both rooted in the same mechanism:

\section*{The Main Cases}

\textbf{1. Topic}
\par\quad 象は鼻が長いです。\\
\par\quad As for elephants — the trunk is long.\\

は sets the frame: everything that follows is understood in relation to "elephants." This is the classic topic function.

\textbf{2. Contrast}
\par\quad コーヒーは飲みますが、紅茶は飲みません。\\
\par\quad I drink coffee, but I don't drink tea.\\

Here は singles out two elements as opposite poles of a contrast. This, too, is singling-out — just with an implicit counterpart instead of a neutral frame.

Both functions share the same core: は explicitly makes something the starting point of the statement, while always leaving something unsaid hovering in the background — the frame in the case of topic, the unstated counterpart in the case of contrast.

\textbf{3. Why が, not は, goes with question words}
\par\quad 誰が来ますか。\\
\par\quad Who's coming?\\

A question word can't be singled out — it's unknown by definition; there's no frame and no contrast to set before the answer is known. That's why が (the identification particle) always goes here, never は.

\section*{An Honest Limit}

With quantity expressions, there's a case that doesn't fit the image quite as cleanly:

\par\quad 20人は集まりました。\\
\par\quad At least 20 people showed up (maybe more).\\

Here は marks something like a threshold with an implicit openness upward — related to the contrast function, but not fully identical to it. This isn't a weakness specific to our image: even in the academic literature, there's no single, fully unified account of every use of は. The singling-out image explains the vast majority of important cases reliably — only at this edge does a small, openly acknowledged fuzziness remain.


\chapter{の — The Thing Particle}

\section*{The Image}

Picture の as an invisible word for "thing" that can be attached to almost anything. の doesn't just connect two nouns (like a possessive) — it can also turn an entire sentence or an adjective into a generic "the X-thing."

These two functions look different at first glance, but they're really the same move: の takes something — a noun, a sentence, an action — and turns it into "the thing belonging to X" or "the thing concerning X."

\section*{The Main Cases}

\textbf{1. Possessive connection}
\par\quad 私の本\\
\par\quad My book (literally: the book belonging to "I")\\

The classic の — connecting two nouns as belonging together.

\textbf{2. Nominalization}
\par\quad 買ったのはこれです。\\
\par\quad What I bought is this. (literally: "the thing of having bought")\\

Here の turns an entire sentence into a "thing," which can then be used just like a noun — as subject, object, or whatever the sentence needs.

\textbf{3. Explanatory sentence ending}
\par\quad どうしたの？\\
\par\quad What's going on? (literally: "what is the thing that...")\\

の at the end of a sentence signals: there's a situation behind this that's being explained or asked about — the same nominalization logic, just at the end of the sentence instead of in the middle.

\section*{A Practical Trick}

の can almost always be mentally replaced with もの (a concrete thing) or こと (a state of affairs) — which makes the nominalization function immediately tangible:

\par\quad 赤いの = 赤いもの\\
\par\quad The red one (literally: the red thing)\\

This also helps with the choice between の and こと as nominalizers — it comes down to a choice between the two nouns の stands in for:

\par\quad 音楽を聞くのが好きです。\\
\par\quad I like listening to music. (の ≈ もの, more concrete, closer to the experience)\\

\par\quad 音楽を聞くことが好きです。\\
\par\quad I like the act of listening to music. (こと, more abstract, more about the fact itself)\\

Both sentences are practically synonymous, but の leans toward the concrete experience, こと toward the abstract state of affairs — fitting with the fact that の historically stands in for もの, and こと simply is こと.

\bigskip\hrule\bigskip

\emph{For the curious:} Historians consider it likely that the nominalization function grew out of the possessive function through an omission: N1のもの ("the thing belonging to N1") probably lost the word もの over time — leaving N1の alone to carry the meaning "the thing belonging to N1" by itself. Both uses would then be, historically, the same construction, just deployed at different points in the sentence.


\chapter{と — The Equivalence Particle}

\section*{The Image}

Picture と as an equals sign: two things are placed on the same level, side by side as equals. と is more than just "with" — it always connects two elements that count as equivalent or matching in some way.

\section*{The Main Cases}

\textbf{1. Partner in a joint action}
\par\quad 友達とカラオケに行きます。\\
\par\quad I'm going to karaoke with a friend.\\

Two equal participants in the same action.

\textbf{2. Result of a change}
\par\quad 氷が水となります。\\
\par\quad Ice becomes water.\\

Here the end state is equated with the starting point — after the change, one \emph{is} the other.

\textbf{3. Quotation}
\par\quad 「こんにちは」と言いました。\\
\par\quad He said, "Hello."\\

What was said and the exact wording are equated — と marks the equivalence between the statement and its precise content.

\textbf{4. Listing}
\par\quad 新聞と雑誌\\
\par\quad Newspaper and magazine\\

The elements stand side by side as equals — the order can be swapped (雑誌と新聞 works just as well), which structurally confirms the equal standing.

\section*{A Sharpened Case That Fits the Image}

With certain verbs like 結婚する (to marry) or 争う (to fight/dispute), the と-partner isn't optional the way it is with mere accompaniment — it's grammatically required, because the action inherently needs two equal participants:

\par\quad 彼女と結婚します。\\
\par\quad I'm marrying her.\\

You can't simply substitute と一緒に ("together with") here — the action of "marrying" only exists between two equal partners. と-partners here aren't mere companions; they're structurally equal co-participants in the action itself.


\chapter{や — The Example Particle}

\section*{The Image}

Picture や as a hand gesturing toward a few things and saying: "things like this, and more of the same." や lists examples without claiming the list is complete.

\section*{Example}

\par\quad 机の上に本や雑誌があります。\\
\par\quad There are books, magazines, and things like that on the desk.\\

The implication: there are probably more things on the desk, but books and magazines are representative examples.

\section*{The Difference from と}

In the と chapter, we saw that と places two or more elements fully and equally side by side — closing the list.

\par\quad 本と雑誌があります。\\
\par\quad There's a book and a magazine. (only these two)\\

\par\quad 本や雑誌があります。\\
\par\quad There are books, magazines, and the like. (probably more)\\

Same sentence structure, but と and や set different expectations: と closes the list, や leaves it open.

\section*{A Small Fuzziness}

In some contexts, や can be used even when all the relevant elements really have been named — the openness becomes more of a linguistic habit than a strict claim about completeness. This isn't an exception to the core idea, just a reminder: や \emph{implies} openness, it doesn't \emph{guarantee} it in every single case.


\chapter{とか — The Example Particle with a Shrug}

\section*{The Image}

とか shares its basic idea with や: naming examples without claiming completeness. But とか goes further — it sounds casual, almost shoulder-shrugging: "stuff like that, who knows what else." It signals not just an open list, but deliberate vagueness.

\section*{The Main Cases}

\textbf{1. Listing, more casual than や}
\par\quad 週末は映画とか見に行きたいな。\\
\par\quad This weekend I'd like to go see a movie or something.\\

Example-like, casual, colloquial.

\textbf{2. Attaching to more than nouns}
\par\quad 忙しいとか疲れたとか、いろいろ理由をつける。\\
\par\quad He comes up with all sorts of excuses — busy, tired, whatever.\\

Unlike や, which only attaches to nouns, とか also works with adjectives and verbs — it's grammatically more flexible.

\textbf{3. Deliberate vagueness}
\par\quad 仕事とか大変なんだ。\\
\par\quad Work and stuff is rough right now...\\

Here the implication is: there are other things the speaker doesn't want to go into detail about. Interestingly, this vagueness effect sometimes shows up even when all the relevant points really have been named (いいとか悪いとか, "good or bad or whatever") — とか gets used to sound deliberately vague, not because anything is actually missing.

\section*{The Difference from や}

Both particles share the core idea (open listing), but とか is more colloquial — best avoided with superiors or in formal speech — and carries this extra tone of casual vagueness that や doesn't have.


\chapter{も — The Accumulation Particle}

\section*{The Image}

Picture も as a plus sign: "this one too." も lifts out an element and adds it to a set of others for whom the same thing already holds. Where は sets up a contrast (this one yes, that one no) and が identifies (exactly this, and only this), も adds (this one too, in addition to what's already known).

\section*{The Main Cases}

\textbf{1. Simple addition}
\par\quad 太郎も学生です。\\
\par\quad Taro is a student too.\\

Taro gets added to a group (other students) for whom the same predicate holds.

\textbf{2. Universal quantifier with negation}
\par\quad 誰も来ませんでした。\\
\par\quad Nobody came.\\

Here も accumulates across \emph{every} possible person — "not this one, not that one, not..." carried all the way through to a complete negation. Same addition, just pushed to its logical end.

\textbf{3. Emphasizing quantity}
\par\quad ビールを10本も飲みました。\\
\par\quad He drank a whopping 10 beers.\\

Here も signals: this amount exceeds expectations — also a form of accumulation, just against an implicit expectation threshold instead of a group of people.

\textbf{4. Extreme example}
\par\quad お茶でも飲みますか。\\
\par\quad Would you like some tea or something?\\

Here a natural example is offered, with the implicit undertone "and anything else along these lines would work too" — the accumulation stays implicit, but points to an open set of alternatives.

\section*{A Combined Case: でも}

でも is built from で (the boundary particle, marking an area) and も (accumulation). That explains its uses directly:

\par\quad 誰でも来られます。\\
\par\quad Anyone can come (no matter who).\\

Here the accumulation isn't just over people (誰も), but over \emph{every conceivable case/domain} — で supplies the domain, も the accumulation across it. でも doesn't need its own separate image — it follows entirely from its two building blocks.


\chapter{から — The Starting-Point Particle}

\section*{The Image}

Picture から as a starting point on a line — the place, time, state, or reason something develops from. から always marks where something begins.

\section*{The Main Cases}

\textbf{1. Spatial starting point}
\par\quad 駅から電車で向かいます。\\
\par\quad I'm heading there by train from the station.\\

\textbf{2. Starting state of a change}
\par\quad 水から氷になりました。\\
\par\quad Water became ice.\\

Here から marks the state \emph{before} the change — the counterpart to と, which (from the と chapter) equates the \emph{end} state with the result (氷が水となる). から looks back at the origin, と marks the result.

\textbf{3. Reason/cause}
\par\quad 忙しいから、行けません。\\
\par\quad Since I'm busy, I can't come.\\

The reason is treated as the starting point of the conclusion — the result follows on from this fact.

\textbf{4. Sequential actions}
\par\quad 食べてから、寝ます。\\
\par\quad After eating, I'll sleep.\\

The first, completed action becomes the starting point for the second.

\section*{The Difference from に with Recipients}

Both に and から can mark the giver with verbs like もらう ("to receive"):

\par\quad 友達に本をもらいました。\\
\par\quad 友達から本をもらいました。\\
\par\quad I received a book from a friend.\\

The difference: に emphasizes the friend as an \textbf{interaction partner} (a point of reference in a relationship), から emphasizes them as a \textbf{pure source} — which is why impersonal sources like institutions almost always take から (組織から奨学金をもらう, "receive a scholarship from an organization"), not に: an organization isn't a relational interaction partner, just a source.


\chapter{まで — The Endpoint Particle}

\section*{The Image}

まで marks the endpoint of a continuous stretch — the point up to which something continues, with that point itself still included. Where から sets the beginning, まで sets the end, and together they span out a stretch.

\section*{The Main Cases}

\textbf{1. Spatial/temporal endpoint}
\par\quad 5ページから10ページまで読みました。\\
\par\quad I read from page 5 to page 10.\\

Page 10 is still included — まで includes the endpoint.

\textbf{2. Inclusive time frame}
\par\quad 15日まで待ってください。\\
\par\quad Please wait until the 15th.\\

The 15th itself is still included.

\section*{The Subtle Difference with までに}

\par\quad 5時まで待ちます。\\
\par\quad I'll wait until 5 o'clock (continuously, the whole time).\\

\par\quad 5時までに終わります。\\
\par\quad It'll be done by 5 o'clock (a deadline, not a continuous stretch).\\

まで alone describes a continuation — the action runs uninterrupted up to that point. Add に (までに), and it becomes a deadline: no longer the stretch itself, but a single point that must be reached by then. The に brings in exactly the point-logic we already know from the に chapter.

\section*{The Bracket with から}

から and まで belong together: から marks where the stretch begins, まで where it ends.

\par\quad 東京から大阪まで新幹線で行きます。\\
\par\quad I'm going from Tokyo to Osaka by shinkansen.\\

Together they form the complete bracket of a fully traversed stretch — matching を, which we know from its own chapter: を marks what gets fully traversed, から...まで marks where that traversal begins and ends.


\chapter{だけ — The Exclusivity Particle}

\section*{The Image}

だけ draws a sharp but neutral boundary around an amount or selection: "exactly this, no more, no less." No regret, no surprise — just a plain statement of the limit.

\section*{Example}

\par\quad 一人だけ来ました。\\
\par\quad Only one person came.\\

Stated neutrally: exactly one, no more.

\section*{The Difference from しか}

Japanese has several words for "only" that differ in tone. だけ is the neutral one:

\par\quad 一人だけ来ました。\\
\par\quad Only one person came. (plain statement)\\

\par\quad 一人しか来ませんでした。\\
\par\quad Only one single person came (and that's too few).\\

しか requires a negative in the sentence and always carries a tone of insufficiency — "that's not enough, too bad." だけ, by contrast, stays neutral: a plain quantity statement with no emotional coloring.

\bigskip\hrule\bigskip

\emph{For the curious:} だけ goes back to the old noun たけ (丈), meaning "measure" or "extent" — as in 背丈 ("height," literally "back-measure"). "Up to the measure of X" fits today's image exactly: だけ draws the boundary right at the measure of what applies.


\chapter{より — The Comparison Starting-Point Particle}

\section*{The Image}

より looks like its own separate "than" particle, but it's actually a specialization of から: instead of a spatial or temporal starting point, より marks the \textbf{starting point of a comparison axis}.

\section*{Example}

\par\quad 東京は大阪より大きいです。\\
\par\quad Tokyo is bigger than Osaka.\\

Literally: "Tokyo is big, measured starting from Osaka." Osaka is the reference point from which the size difference is measured — just as から marks a starting point for movement, より marks a starting point for comparison.

\section*{Why This Isn't Just a Metaphor}

Unlike most particle images, the connection between より and から isn't just a conceptual bridge — it's directly attested historically: in Classical Japanese, より was actually the primary starting-point particle, from which today's comparison function developed later as a special case. The old starting-point meaning survives to this day in formal, written language:

\par\quad 5ページより先を読んでください。\\
\par\quad Please continue reading from page 5 onward.\\

Here より functions exactly like から — just in a more formal register.

\section*{A Special Construction}

\par\quad これより他に方法がない。\\
\par\quad There's no other way besides this one.\\

Here too, より is a starting point: "starting from this option, there's no other" — the same idea, just combined with negation to form an exclusion construction.


\chapter{し — The Accumulating Reason Particle}

\section*{The Image}

し works like も, but at the sentence level: while も adds a noun to a set of others ("this one too"), し adds a fact or reason to others that also apply — "and on top of that, this too."

\section*{The Main Cases}

\textbf{1. Multiple reasons listed}
\par\quad 物価も高いし、電車も人が多いし、それに暑い。\\
\par\quad Prices are high, the trains are crowded, and on top of that it's hot.\\

Each reason gets attached with し — the list stays open, more reasons could follow.

\textbf{2. A single reason, with others implied}
\par\quad 疲れたし、もう寝ます。\\
\par\quad I'm tired, so I'm going to sleep now.\\

Even when only one reason is given, し often implies: there are probably more reasons that just aren't being said. This is what sets し apart from から, which states a reason without implying a whole collection behind it — から is the plain, single justification; し hints at a larger set.

\section*{Why し Sounds Less Assertive Than から}

Because し always hints at an open set of contributing reasons, a justification made with し often sounds less categorical or final than one made with から — it comes across more like "among other things, because..." rather than "precisely because..."


\chapter{けれど — The Concession Particle}

\section*{The Image}

けれど connects two parts of a sentence where the second goes against the expectation set up by the first: "Even though this is true, that's still the case." It's the everyday, spoken counterpart of "although" or "but."

\section*{Example}

\par\quad 高いけれど、買いました。\\
\par\quad It was expensive, but I bought it anyway.\\

The first part (expensive) sets up a certain expectation (don't buy it) — the second part contradicts that expectation.

\section*{Kinship with のに}

We already know this function from the の chapter: のに works on the same pattern — one clause sets up an expectation that the following clause doesn't fulfill. けれど is, in a sense, the standalone conjunction version of the same concessive idea, while のに ties directly into の's nominalization function.

\section*{Register}

けれど appears in various shortened forms — けれども (most formal), けれど, けど (most casual) — depending on how relaxed the speech situation is. The meaning stays the same across all forms.


\chapter{ながら — The Coexistence Particle}

\section*{The Image}

ながら says: two things hold true at the same time, without canceling each other out. Most familiar for simultaneous actions, but the same core idea carries three further cases.

\section*{The Main Cases}

\textbf{1. Simultaneous actions}
\par\quad 音楽を聴きながら勉強します。\\
\par\quad I study while listening to music.\\

Two actions running alongside each other.

\textbf{2. Concession — two coexisting states that seem incompatible}
\par\quad 貧しいながらも幸せな家庭です。\\
\par\quad A family that is poor and yet happy.\\

Here "poor" and "happy" hold true at the same time, even though one wouldn't necessarily expect them together. The "although" meaning arises naturally out of the coexistence: when two simultaneously true things stand in tension, that reads automatically as a contrast.

\textbf{3. Persisting state}
\par\quad 昔ながらのたたずまい。\\
\par\quad An appearance unchanged since the old days.\\

The old state persists all the way into the present — nothing has changed.

\textbf{4. Totality}
\par\quad 三度ながら失敗しました。\\
\par\quad I failed all three times.\\

Multiple occurrences are treated as one combined whole — also a form of coexistence: the three times "exist together" as a single unit.

\section*{A Rule of Thumb}

With action verbs (音楽を聴きながら), ながら usually reads as plain simultaneity. With adjectives or state words (貧しいながら), it tends to read as concessive. This isn't a strict rule, but a useful first orientation.


\chapter{か — The Uncertainty Particle}

\section*{The Image}

か marks that a piece of information isn't settled — it lies outside the speaker's control or certainty. The most familiar use is the question, but that's just the most visible case of a much broader underlying idea.

\section*{The Main Cases}

\textbf{1. Question}
\par\quad あなたも行きますか。\\
\par\quad Are you going too?\\

The information sits with the listener — the speaker doesn't know.

\textbf{2. Indefiniteness with question words}
\par\quad 誰かが来ました。\\
\par\quad Someone came.\\

Here か attaches to a question word and deliberately leaves it open who exactly is meant — turning it into an indefinite reference rather than a real question.

\textbf{3. Alternative}
\par\quad コーヒーか紅茶か選んでください。\\
\par\quad Please choose coffee or tea.\\

Which of the two options applies stays open — か marks exactly this undecided state.

\textbf{4. Self-questioning}
\par\quad どうしようか。\\
\par\quad What should I do?\\

Interestingly, here the speaker treats even their own not-yet-made decision as uncertain — か shows that even the speaker doesn't yet know the answer.

\textbf{5. Surprise at new information}
\par\quad だれかと思ったら、君だったのか。\\
\par\quad Oh, so it was you!\\

A fact just learned that hasn't fully been "processed" or integrated yet — also a form of uncertainty, just tied to the moment of surprise.

\section*{A Bigger Pattern}

か fits into a series of particles we already know, all concerned with the degree of certainty in a statement: と (と chapter) sets a complete, fixed list. や (や chapter) deliberately leaves open whether more belongs to it. とか (とか chapter) adds a deliberate vagueness on top of that. か goes furthest of all: here it's not just the completeness of a list left open, but the information itself.


\chapter{ね — The Particle of Shared Certainty}

\section*{The Image}

ね says: "We both know this already, right?" It marks a piece of information as shared property between speaker and listener, and seeks confirmation or agreement about it. ね draws the participants closer together — it invites the listener to resonate along.

\section*{Example}

\par\quad いい天気ですね。\\
\par\quad Nice weather, isn't it?\\

The speaker assumes the listener sees and feels the same thing, and seeks confirmation of this shared perception.

\par\quad もうすぐ夏休みですね。\\
\par\quad Summer break is coming up soon, isn't it?\\

Again, a fact both conversation partners already know is jointly affirmed — not to deliver information, but to build closeness and agreement.


\chapter{よ — The Particle of Exclusive Speaker Knowledge}

\section*{The Image}

よ says: "You don't know this yet, but I'm telling you now." Unlike ね, which confirms shared knowledge, よ shares information that (in the speaker's assumption) belongs to them alone and is new to the listener. よ often also carries emphasis — the speaker insists on their point of view.

\section*{The Main Cases}

\textbf{1. Sharing new information}
\par\quad わたしは行きませんよ。\\
\par\quad I'm not going, I'm telling you!\\

The speaker shares a fact the listener doesn't yet know.

\textbf{2. Drawing attention, adding emphasis}
\par\quad 危ないよ！\\
\par\quad Watch out!\\

Here よ heightens the urgency — the speaker wants to make sure the warning lands.

\section*{The Difference from ね}

In the ね chapter, we saw: ね assumes shared knowledge and seeks confirmation of it. よ flips that around — it shares something the speaker believes they alone know, and the listener doesn't.

\par\quad 明日、雨ですよ。\\
\par\quad It's going to rain tomorrow — you didn't know that.\\

\par\quad 明日、雨ですね。\\
\par\quad It's going to rain tomorrow, isn't it? (we both already know that)\\

Both particles answer the same silent question — "where does this knowledge sit, with me alone or with both of us?" — just with opposite answers.


\chapter{かい — The Casual Question Variant}

\section*{The Image}

かい is built from か (the uncertainty particle we already know) and い — two question-building blocks combined into a particularly familiar, colloquial question form. かい doesn't carry any meaning beyond か on its own; it's essentially the same uncertainty marker, just with an added tone of familiarity or casualness.

\section*{Example}

\par\quad 元気かい？\\
\par\quad How's it going?\\

Same question as 元気ですか, but noticeably more casual and familiar in tone — typical among close friends or in very informal speech.

\section*{A Note on Usage}

かい is often described as purely masculine speech. That does match today's perception fairly well — in pop culture (e.g. yakuza films) it reads as strongly masculine. Historically, though, かい wasn't exclusively restricted to male speakers; the strict association is more of a modern stereotype than a hard-and-fast rule. In practice today, かい tends to sound masculine-casual in most contexts and should be used with that in mind.


\chapter{だい — The Casual Wh-Question Variant}

\section*{The Image}

だい is the counterpart to かい from the previous chapter — the same idea of a soft, familiar question variant, just for a different kind of question. While かい goes with plain yes/no questions, だい goes with questions containing a question word (who, what, where, when...).

\section*{Example}

\par\quad どこに行くんだい？\\
\par\quad Where are you off to, then?\\

\par\quad それ、何だい？\\
\par\quad What's that, then?\\

The same question could be asked neutrally with か (どこに行くんですか) — だい makes it noticeably more casual and familiar, typical of older, informal speech.

\section*{The Division of Labor with かい}

Both forms — かい for yes/no questions, だい for wh-questions — work on the same principle: an added vowel softens a direct question and makes it more familiar. かい attaches to か, だい to だ.

\section*{A Note on Usage}

だい tends to sound predominantly masculine-casual in practice, with the well-known exception that it's also traditionally used by older women. As with かい: use with an awareness of its distinct register.


\chapter{Verb Stem Forms — The Toolkit Behind the Forms}

\section*{A Different Kind of Chapter}

So far, every chapter has been about a particle — a choice a speaker makes among several images. Stem forms are something else: they're the building blocks that almost all verb forms are built from — negation, polite form, command form, conditional clauses, and much more. A verb isn't "used" in a stem form; it's built onward through it into another form.

\section*{The Practical Core: Ending Vowels}

For verbs with a consonant stem ("u-verbs," e.g. 書く), you can identify the stem form directly from its ending vowel. In Japanese vowel order (a–i–u–e–o):

| Vowel | Stem Form | Example (書く) | Leads to |
|---|---|---|---|
| -a | Mizen-kei | 書か | ない-form, せる (causative), れる (passive) |
| -i | Ren'yō-kei | 書き | ます-form, たい-form |
| -u | Shūshi-kei | 書く | dictionary form |
| -e | Katei-kei/Meirei-kei | 書け | ば-form (conditional), command form |
| -o | Mizen-kei 2/Suiryō-kei | 書こ | よう-form (volitional/presumptive, +う) |

音便形 (the "sound-change form") has no ending vowel of its own — it arises as a sound shift of the ren'yō-kei before て/た (書き+て → 書いて).

\section*{Other Verb Groups}

For verbs whose stem ends in a consonant (e.g. 書く, 話す, 読む — often called "u-verbs" because the dictionary form ends in -u), the vowel shift is clearly visible, since the stem itself carries the vowel. Other verb groups work differently:

\textbf{Verbs with a stem ending in -eru/-iru (e.g. 食べる, 見る)} — the stem stays unchanged; only the ending changes:

| Stem Form | 食べる |
|---|---|
| Mizen-kei | 食べ (+ない, +させる, +られる) |
| Ren'yō-kei | 食べ (+ます, +たい) |
| Shūshi-kei | 食べる |
| Katei-kei/Meirei-kei | 食べれ (+ば) / 食べろ |
| Mizen-kei 2/Suiryō-kei | 食べよ (+う) |

\textbf{The two irregular verbs} — する and 来る don't follow either pattern consistently and have to be learned individually:

| Stem Form | する | 来る |
|---|---|---|
| Mizen-kei | し (+ない), さ (+せる), さ (+れる) | 来 (こ) (+ない) |
| Ren'yō-kei | し (+ます, +たい) | 来 (き) (+ます, +たい) |
| Shūshi-kei | する | 来る (くる) |
| Katei-kei/Meirei-kei | すれ (+ば) / しろ | 来れ (くれ) (+ば) / 来い (こい) |
| Mizen-kei 2/Suiryō-kei | し (+よう) | 来 (こ) (+よう) |

Seven traditional stem forms (really six, plus one later split) make up this foundation. For some of them, every form built from that stem shares the same underlying idea; for others, only some do — and that's exactly what makes them interesting.

\section*{The Building Blocks at a Glance}

\textbf{Ren'yō-kei (the "connective" form)} — works like の: it turns the verb into a "thing" that can attach to other things (the ます-form, the たい-form, or combinations with other words). This nominalization idea carries consistently through every form built from this stem.

\textbf{Shūshi-kei (the "terminal" form)} — the dictionary form, used to fully close out a sentence.

\textbf{Mizen-kei (the "not-yet" form)} — this is where it gets interesting: the ない-form (negation — something is "not yet" the case) and the presumptive form fit this idea. The causative (せる) and passive (れる), also built from this same stem, don't — they attach to it purely for phonetic reasons (a connecting vowel is needed), not because of any shared "not yet" idea.

\textbf{Katei-kei (the "hypothetical" form)} — similarly mixed: the ば-conditional carries the "if" idea, but other forms built from the same stem don't carry it in the same way.

\textbf{Ishikei — two forms wrongly treated as one.} What was traditionally treated as a single stem form is actually two independent building blocks: the command form (meirei-kei) and the presumptive/volitional form (suiryō-kei, the basis for the よう-form). They don't even share the same stem vowel.

\textbf{On'binkei (the "sound-change" form)} — describes a purely phonetic process: the て-form and た-form arise from the same pronunciation-easing mechanism, carrying no meaning of their own.

\section*{What This Means for the Chapters Ahead}

The chapters that follow deal with the concrete forms built from these building blocks — the ます-form, the たい-form, the ない-form, the causative, the passive, and so on. There, what matters is each form's own meaning, not which building block it's technically constructed from. This chapter is only background knowledge — for understanding why certain forms are morphologically related even though they have nothing to do with each other in meaning.


\chapter{The ない-Form — Negation}

\section*{The Image}

ない marks absence — that something is not the case, doesn't exist, or doesn't happen. Whether attached to a verb or used as a standalone word, the core idea stays the same: something is missing.

\section*{The Main Cases}

\textbf{1. Verb negation}
\par\quad 食べない。\\
\par\quad I'm not eating.\\

The action is absent — it doesn't take place.

\textbf{2. As the negation of ある: "doesn't exist"}
\par\quad お金がない。\\
\par\quad I don't have money. (literally: money doesn't exist.)\\

Where ある ("to exist, to be present") asserts a thing's existence, ない negates it — the same absence idea as with verb negation, just applied to a thing instead of an action.

\textbf{3. Advice}
\par\quad 早く寝た方がいいよ。\\
\par\quad You'd better go to bed early.\\

The negation form is embedded here in the comparison construction (方がいい, "it would be better..."), implicitly pointing at the absence of the worse alternative.

\textbf{4. About to not do something}
\par\quad 食べないところだ。\\
\par\quad I'm about to not eat (I'm not eating right now).\\

Again, the absence of the action becomes the central point of the statement.

\section*{A Small Detail}

The negation of verbs (読まない, "not read") and the negation of existence (お金がない) originally have different roots — historically they aren't the same word. Over time they came to share the same conjugation pattern. Even so, both still carry the same core idea today: something isn't there.


\chapter{The たい-Form — Subjective Desire}

\section*{The Image}

たい expresses a desire known from the inside — one's own. It's the form used to say what you yourself want, not what you observe in others.

\section*{Example}

\par\quad 日本に行きたいです。\\
\par\quad I want to go to Japan.\\

A desire expressed from one's own, immediate perspective.

\section*{The Limit: Only for Yourself (and, in Conversation, the Listener)}

たい is normally only used for oneself, or in questions for the conversation partner:

\par\quad 何が食べたいですか。\\
\par\quad What do you want to eat?\\

For a third person, however, たい doesn't quite fit — you can't see directly into someone else's desire from the outside, only observe their behavior:

\par\quad 彼は日本に行きたがっています。\\
\par\quad He seems to want to go to Japan. (observed, not known directly)\\

This calls for the form たがる — it describes a desire as it appears from the outside, not as it's experienced from the inside. This distinction — what the speaker knows directly, versus what they can only observe from the outside — is a recurring principle in Japanese: the same logic shows up with adjectives describing feelings (e.g. 痛い, "painful"), which also take the がる-form for a third person (痛がる, "to show signs of pain").


\chapter{The ます-Form — The Polite Form}

\section*{The Image}

ます turns a neutral statement into a polite one — the form used to show respect and appropriate distance toward a listener, without coming across as cold or impersonal. It's the standard level of politeness used in most everyday situations with strangers, colleagues, or people deserving respect.

\section*{Example}

\par\quad 毎日、日本語を勉強します。\\
\par\quad I study Japanese every day.\\

The same statement without ます (勉強する) means the same thing, but sounds direct and casual — fine among friends, but not in a formal context.

\section*{Why This Form Expresses Politeness}

ます is generally traced back to an old, humble verb, with which one originally signaled a degree of self-effacement toward the person a certain action concerned — though the precise origin is not settled in the research. On this reading, the modern, broader politeness toward the listener as such developed over time out of this original humility toward the person a given action was directed at — regardless of what the sentence is actually about.


\chapter{よう — Two Forms That Sound Alike}

\section*{A Warning Up Front}

よう looks like a single form at first glance. In reality, it hides two completely independent building blocks that just happen to sound the same — much like the English word "bark," which can mean the outer layer of a tree or the sound a dog makes, with no connection between the two meanings at all. Treating both よう as the same thing quickly leads to confusion. That's why they're deliberately kept separate here.

\section*{The First よう: Will and Presumption}

This form comes directly from the verb itself (行こう, 食べよう) and expresses an expression of will or a presumption:

\par\quad 一緒に行こう。\\
\par\quad Let's go together!\\

\par\quad もうすぐ雨が降ろう。\\
\par\quad It'll probably rain soon.\\

This よう doesn't have to sit at the end of the sentence — it can also appear in the middle, for instance before と:

\par\quad 早く行こうと思っています。\\
\par\quad I'm thinking of going soon.\\

\section*{The Second よう: Comparison, Likelihood, Purpose}

This form goes back to the noun 様 ("manner, way") and mostly appears as ようだ or ように:

\par\quad 雨が降るようだ。\\
\par\quad It looks like it's going to rain.\\

\par\quad 猫のように寝ている。\\
\par\quad He's sleeping like a cat.\\

\par\quad 忘れないように書いておきます。\\
\par\quad I'm writing it down so I won't forget.\\

Here it's not about will at all, but about resemblance, appearance, or purpose — content-wise this has nothing to do with the first よう, even though it sounds identical.

\section*{Where the Two Collide}

One case makes the risk of confusion especially clear:

\par\quad どうしようもない。\\
\par\quad There's nothing that can be done about it.\\

Here よう sounds like the volitional よう (しよう, "I will do"), but actually belongs to the second family — 様 in the sense of "way, method" ("there's no way/method to do it"). Cases like this show exactly why it's worth treating the two よう as separate from the start, even though they sound the same.

\bigskip\hrule\bigskip

\emph{For the curious:} Historically, the two よう are unrelated in origin — one grows out of the verb's own volitional stem, the other out of the noun 様. Their identical sound today is coincidence, not shared history.


\chapter{The そう-Form — Impression and Hearsay}

\section*{The Image}

そう marks that a statement isn't based on direct personal knowledge, but on an outward impression or on something the speaker heard. There are two uses that look very similar on the surface, but attach at different points in the sentence — and that's exactly how you can tell them apart.

\section*{Impression (attaches directly to the stem)}

\par\quad このケーキ、美味しそうです。\\
\par\quad This cake looks delicious.\\

Here そう attaches directly to the word stem (美味し-, without the -い) — it describes an immediate, visual or sensory impression. The speaker hasn't tried the cake yet, but it looks that way.

\section*{Hearsay (attaches to the full form)}

\par\quad 天気予報によると、明日は雨だそうです。\\
\par\quad According to the weather forecast, it's supposed to rain tomorrow.\\

Here そう attaches after the complete form (雨だ, not just 雨) — it reports what the speaker heard from someone else, not what they perceive directly.

\section*{The Key Difference}

Both uses look similar, but the attachment point reveals the meaning:

\par\quad 面白そうです。 → Impression: looks interesting (stem + そう)\\
\par\quad 面白いそうです。 → Hearsay: is said to be interesting (full form + そう)\\

A single い decides between "I get the impression" and "I heard that." Once you know this attachment rule, the two uses are easy to tell apart with confidence.


\chapter{The た-Form — The Past Tense}

\section*{The Image}

た marks that something has already happened — the past. Behind it lies always the same idea: from the speaker's point of view, at the moment of speaking, an event already lies behind them.

\section*{Only Two Time Frames}

English distinguishes past, present, and future as three separate tenses. Japanese grammar knows only one boundary: \textbf{past} versus \textbf{non-past}. The dictionary form of a verb (行く, "go") covers both present and future:

\par\quad 明日、東京に行く。\\
\par\quad I'm going to Tokyo tomorrow. (future, dictionary form)\\

\par\quad いつも歩いて行く。\\
\par\quad I always go on foot. (present/habitual, dictionary form)\\

た, on the other hand, marks only that something has already happened:

\par\quad 昨日、映画を見た。\\
\par\quad I watched a movie yesterday.\\

So the only grammatical boundary runs between "already happened" and "not yet happened" — whether something is happening right now or will happen in the future isn't distinguished by た at all, and has to be made clear, if needed, through time words in the sentence (明日, いつも, and so on).

\section*{Even a Just-Reached State Is Already the Past}

Any state that has just been reached is, by definition, already a moment in the past the instant you notice it:

\par\quad あ、財布が落ちた！\\
\par\quad Oh, my wallet fell!\\

Even though this exclamation happens at the very moment of noticing, the falling itself already lies a tiny moment in the past — the wallet is already on the floor by the time you notice it. た captures exactly this point: the action is already complete, its result has just become visible.

\par\quad 結婚しています。\\
\par\quad I'm married.\\

Here too: a state that arose from a completed action and continues to hold — the same past-tense logic, just applied to an ongoing state.

\textbf{An everyday example:} When you've just understood something, you say 分かりました, not 分かります — even if the understanding happened only a fraction of a second ago. At the moment you say this sentence, the actual act of understanding already lies behind you — you're already in the state of having understood, and that state, like any reached state, is already the past.

\bigskip\hrule\bigskip

\emph{For the curious:} た goes back historically to a combination of て and ある — literally "an action has been carried out and the resulting state exists." This connection to ある explains why た has been so closely tied to the idea of an existing, reached state from the very start.


\chapter{The て-Form — Chaining}

\section*{The Image}

て connects two actions or states in a way that carries a sense of sequence or connection: first this, then (or because of this) that. Unlike と (と chapter), which places two elements truly on equal footing and interchangeably, て always carries a trace of sequence or connection.

\section*{The Main Cases}

\textbf{1. Temporal sequence}
\par\quad 朝ご飯を食べて、学校に行きます。\\
\par\quad I eat breakfast and (then) go to school.\\

First one, then the other — the order is fixed, unlike a true listing.

\textbf{2. Reason}
\par\quad 熱があって、休みました。\\
\par\quad Because I had a fever, I rested.\\

Here the sequence shifts from purely temporal to causal — the same chaining principle, just with emphasis on causation rather than plain order.

\textbf{3. Accompanying action}
\par\quad 音楽を聴いて、勉強しています。\\
\par\quad I'm studying while listening to music.\\

Two things happening closely tied together — again a chain, not two independent, interchangeable facts.

\textbf{4. Polite request}
\par\quad ちょっと待ってください。\\
\par\quad Please wait a moment.\\

てください builds directly on the same chaining principle: the action is meant to follow on from the request.

\section*{The Difference from と}

と (と chapter) places two elements truly on equal footing — in a listing with と, the order can be swapped freely without changing the meaning. With て, there's almost always a direction built in: first this, then that, or because of this, that. That's why て suits sentences meant to express a sequence or connection, while と is better suited to plain, neutral listing.


\chapter{The ば-Form — The Conditional}

\section*{The Image}

ば marks a condition: "if X, then Y" — a general, often rule-like link between two circumstances.

\section*{Examples}

\par\quad 春になれば、桜が咲きます。\\
\par\quad When spring comes, the cherry blossoms bloom.\\

A condition that sounds like a general rule: as soon as one condition holds, the other follows almost automatically.

\par\quad ボタンを押せば、ドアが開きます。\\
\par\quad If you press the button, the door opens.\\

\par\quad 安ければ、買います。\\
\par\quad If it's cheap, I'll buy it.\\

\section*{The Character of ば}

ば often sounds more general and rule-like than other conditional forms — fitting for proverbs, natural laws, fixed rules, or instructions, where a condition leads to a consequence in an almost mechanical way.


\chapter{なら(ば) — Condition and Topic}

\section*{The Image}

なら works at its core like ば from the previous chapter, except the condition attaches not to a verb but to a statement ("it is so"). Out of this arise two closely related uses: a condition, and — somewhat less expected — a kind of topic marker.

\section*{Condition}

\par\quad 雨なら、行きません。\\
\par\quad If it's raining, I won't go.\\

The same conditional logic as ば, just applied to an already-established situation ("it's raining") rather than an unfolding event.

\section*{Topic}

\par\quad 母なら、間もなく帰ります。\\
\par\quad As for Mother, she'll be back soon.\\

Here なら behaves almost like は — it singles out a topic and sets up a statement about it. This second use isn't a random side meaning; it's already latent in the conditional logic itself: "if we're talking about X..." is, in the end, also a kind of condition, just applied to a conversation topic instead of an event.


\chapter{The る-Form — The Potential Form}

\section*{The Image}

Some verbs have their own, direct way of expressing "can" — not as an add-on construction, but as a built-in potential form of the verb itself.

\section*{Examples}

\par\quad 漢字が読める。\\
\par\quad I can read kanji.\\

読める is the standalone potential form of 読む ("to read") — not a passive, not a politeness form, but its own verb that directly means "can read."

\par\quad 泳げる。\\
\par\quad I can swim.\\

Same pattern: 泳ぐ ("to swim") becomes its own potential form, 泳げる.

\section*{An Important Distinction}

This potential form is not the same as れる/られる (covered in the next chapter) — both eventually lead to the meaning "can," but they're two different, independently formed constructions. The る-form comes directly from the verb stem itself and describes a property residing in the object or the action itself — "this is readable," not "I am granted permission to read." Treating both forms as identical confuses two grammatically separate paths to the same meaning territory.


\chapter{れる/られる — Four Meanings, One Origin}

\section*{The Image}

れる (or られる, depending on the verb group) traces back to the idea that something "arises by itself" or "happens on its own" — without anyone actively bringing it about. Out of this single core idea, four seemingly very different uses develop.

\section*{The Four Uses}

\textbf{1. Passive}
\par\quad 先生に褒められました。\\
\par\quad I was praised by the teacher.\\

The action "arises" at the subject, without the subject itself acting — it happens to them.

\textbf{2. Spontaneous feeling}
\par\quad 昔のことが思い出される。\\
\par\quad The old days come to mind (on their own).\\

Here the core idea shows most directly: a feeling or thought arises by itself, with no conscious intention behind it.

\textbf{3. Possibility}
\par\quad この水は飲まれる。\\
\par\quad This water is drinkable.\\

A possibility "arises" — something becomes feasible without anyone actively doing anything to make it so.

\textbf{4. Honorific}
\par\quad 先生が話される。\\
\par\quad The teacher speaks. (respectful)\\

For a respected person, the action is presented as if it "arises on its own," rather than being attributed directly and immediately to the person — an indirect, polite kind of distancing.

\section*{Why All Four Belong Together}

What connects passive, spontaneous feeling, possibility, and honorific use is the same underlying idea: the action isn't presented as actively originating from the subject, but as something that happens or arises. With the honorific use, this distancing is even deliberately exploited as a politeness device — avoiding attributing the action too directly to the honored person by having it grammatically "happen on its own."


\chapter{せる/させる — Making Someone Act}

\section*{The Image}

せる (or させる) turns a verb into an action where someone causes another person to do something — ranging from outright coercion to mere permission, depending on how the construction is used.

\section*{Example}

\par\quad 子供に野菜を食べさせます。\\
\par\quad I'm making the child eat vegetables. / I'm letting the child eat vegetables.\\

Between coercion and permission lies a whole range — the sentence alone allows for either reading; tone and context decide.

\section*{A Spectrum, Not a Fixed Point}

In the を chapter, we saw that particle choice with the causative makes this range visible:

\par\quad 子供を勉強させます。 → with を: more like coercion — the child is "pushed all the way through" studying.\\
\par\quad 子供に勉強させます。 → with に: more like permission — the child is allowed to study.\\

Causative and permission aren't two separate meanings, then, but two ends of the same spectrum — strong coercion on one side, loose permission on the other. Which end is meant is often signaled by the particle choice itself.

\section*{A Polite Extension}

\par\quad 帰らせていただきます。\\
\par\quad I'd like to (with your permission) take my leave.\\

Here the causative form combines with いただく — literally "I receive the permission to [leave]." This expresses particular humility: the speaker humbly asks permission for their own action instead of simply carrying it out.


\chapter{欲しい / 欲しがる — Wanting a Thing}

\section*{The Image}

欲しい expresses wanting a thing — the counterpart to たい (たい-form chapter), which applies to an action, while 欲しい applies to a thing.

\section*{Example}

\par\quad 新しい靴が欲しいです。\\
\par\quad I want new shoes.\\

\section*{The Same Limit as with たい}

Just like たい, 欲しい is normally only used for oneself, or in conversation for the listener:

\par\quad 何が欲しいですか。\\
\par\quad What do you want?\\

For a third person, we again need the がる-form we already know from たがる:

\par\quad 彼女は新しい靴が欲しがっています。\\
\par\quad She seems to want new shoes. (observed from the outside)\\

欲しがる describes a desire visible from the outside in another person, not a desire known directly as one's own — the same principle as たい/たがる, just applied to things instead of actions.


\chapter{みたいだ — The Colloquial Variant of ようだ}

\section*{The Image}

みたいだ works like ようだ from the よう chapter — comparison, likelihood, appearance — just more casual and colloquial in tone.

\section*{Examples}

\par\quad 雨が降るみたいだ。\\
\par\quad Looks like it's going to rain.\\

\par\quad 猫みたいに寝ている。\\
\par\quad He's sleeping like a cat.\\

The same meanings as ようだ/ように, just more common in casual, everyday speech.

\section*{Where the Casual Tone Comes From}

みたいだ actually developed from 見たようだ ("it looks as if seen") and gradually got worn down phonetically into みたいだ — not just a similar-sounding word, but a direct, colloquial evolution of ようだ itself. Its casual character follows directly from this origin: an everyday, worn-smooth pronunciation of the same underlying idea.


\chapter{だ / です / である — Stating a Fact in Three Registers}

\section*{The Image}

だ, です, and である are all, at their core, the same thing: they state that something is the case — "X is Y." The difference between the three lies not in meaning, but purely in the level of politeness.

\section*{The Three Forms}

\par\quad これは本だ。\\
\par\quad This is a book. (direct, casual)\\

\par\quad これは本です。\\
\par\quad This is a book. (polite, standard)\\

\par\quad これは本である。\\
\par\quad This is a book. (formal, written register — essays, reports, official texts)\\

All three say exactly the same thing — only the tone differs.

\section*{Why だ Sounds More Direct Than It First Appears}

だ has a reputation for sounding especially blunt or abrupt, almost like a full stop after the statement: "That's how it is. Done." One tempting but mistaken explanation: だ sounds harsher because it can't be negated. That's not actually true — だ can absolutely be negated (これは本ではない/じゃない), just like です (これは本ではありません). Both need a separate replacement word for this (ではない/じゃない); neither has its own built-in negative form.

The real reason for だ's blunt sound lies elsewhere: in natural speech, だ is very often simply dropped to soften a sentence (本。 instead of 本だ。) — です, by contrast, almost never gets dropped, since it's the standard polite sentence-ender. This droppability, combined with だ's generally blunt, casual tone, is what creates the impression that だ is "harsher" — not the negation.


\chapter{Honorific Speech — Who Gets Raised, Who Lowers Themselves}

\section*{The Image}

Japanese politeness doesn't work through a single level of formality — it works through the question: who gets elevated, and who lowers themselves? There are three fundamentally different directions of movement that make up honorific speech.

\section*{The Three Directions}

\textbf{1. Elevating the other person (尊敬語, sonkeigo)}
\par\quad 先生がいらっしゃいます。\\
\par\quad The teacher is coming. (respectful, because it concerns a higher-status person)\\

The action of the listener or a respected person is grammatically "elevated" — special verb forms replace the neutral ones (いらっしゃる instead of 来る/行く/いる).

\textbf{2. Lowering oneself (謙譲語, kenjōgo)}
\par\quad 先生のところに伺います。\\
\par\quad I will go see the teacher. (humble, because the speaker's own action is played down relative to the respected person)\\

Here the speaker lowers themselves grammatically, indirectly elevating the other person — the same underlying idea as sonkeigo, just approached from the opposite direction.

\textbf{3. General politeness (丁寧語, teineigo)}
\par\quad 明日、行きます。\\
\par\quad I'm going tomorrow. (simply polite, without elevating or lowering anyone in particular)\\

We already know this from the ます-form chapter — a polite tone toward the listener, regardless of who appears in the action itself.

\section*{A Subtlety in Self-Lowering}

Humble speech (kenjōgo) comes in two varieties that differ in one important respect:

\textbf{Directed humility} — the action concerns a specific person, who is thereby elevated:
\par\quad 先生に伺います。\\
\par\quad I will ask the teacher. (lowering oneself explicitly before the teacher)\\

\textbf{Undirected humility} — a generally formal, humble speech style, without addressing or elevating any specific person:
\par\quad 明日、参ります。\\
\par\quad I will come tomorrow. (a formal register, without elevating any particular listener)\\

The difference can be pictured this way: the first variant is like a bow to a specific person; the second is more like a suit you wear to any formal occasion — regardless of who's actually there.


\chapter{Adjectives — Two Very Different Families}

\section*{The Image}

Japanese has two entirely different kinds of "adjectives" that look similar on the surface but work differently grammatically: い-adjectives, which conjugate on their own like verbs, and な-adjectives, which behave more like nouns.

\section*{い-Adjectives}

\par\quad 高い山\\
\par\quad a tall mountain\\

\par\quad この山は高い。\\
\par\quad This mountain is tall.\\

い-adjectives can stand directly before a noun (attributive) or directly end a sentence (predicative), with no copula (だ/です) needed in between. They carry their own conjugation (高かった, "was tall"; 高くない, "is not tall").

\section*{な-Adjectives}

\par\quad 静かな部屋\\
\par\quad a quiet room\\

\par\quad この部屋は静かです。\\
\par\quad This room is quiet.\\

な-adjectives need an attached な before a noun, and the copula だ/です to end a sentence — exactly like an ordinary noun (本の部屋 would have the same structure as 静かな部屋, just with の instead of な).

\section*{Why な-Adjectives Are Really Nouns}

This resemblance to ordinary nouns is no coincidence. There's a well-established position in Japanese grammatical research holding that な-adjectives aren't a separate word class at all, but simply nouns combined with だ — 静かな would then really be the noun 静か plus a form of だ, just like an ordinary noun. This view is supported by the conjugation pattern itself: the inflected forms of a な-adjective (だろ, だっ/で/に, だ, な, なら) are identical to the conjugation of the copula だ itself, which we already know from the だ/です/である chapter.

In practice, this means: it's easiest to think of な-adjectives as nouns with an especially close relationship to だ — not as their own verb-like word class the way い-adjectives are.


\chapter{Nouns and Verbal Nouns — When a Noun Becomes a Verb}

\section*{The Image}

Many Japanese nouns — mostly of Chinese origin — can combine directly with する ("to do") to form a verb. Instead of learning a separate verb for every action, it's often enough to say: noun + する.

\section*{Example}

\par\quad 運動をします。\\
\par\quad I exercise.\\

\par\quad 運動します。\\
\par\quad I exercise. (same meaning)\\

を is optional here — as long as there's no other object in the sentence, it can be dropped without changing the meaning.

This optionality only applies to the bare noun, though. As soon as an adjective modifies the noun, を becomes required:

\par\quad 難しい運動をします。\\
\par\quad I do difficult exercise.\\

\par\quad ×難しい運動します。\\
\par\quad (not correct)\\

So once an adjective is placed in front, を must stay — only with the plain, standalone noun can it be dropped.

\section*{An Extension: Perception and Impression}

\par\quad いいにおいがする。\\
\par\quad It smells good.\\

\par\quad 変な音がする。\\
\par\quad It sounds strange.\\

Here する works differently — not as "carrying out an action," but as "a perception arises." Even so, the basic principle holds: a noun (におい, "smell") is turned into a complete expression through する.

\section*{A Productive, Living Construction}

This noun+する construction remains highly productive in modern Japanese — it even gets applied to nouns that weren't traditionally meant for it:

\par\quad 青春する。\\
\par\quad To fully live out one's youth. (literally: "to do youth")\\

Creative formations like this show just how flexible this simple pattern is: almost any suitable noun can be turned into its own action with する.


\chapter{Adverbs — Direct, Derived, or Disguised as Nouns}

\section*{The Image}

Japanese adverbs come in three forms: standalone adverbs, adverbs derived from adjectives, and a third, less obvious group — nouns that take on adverbial functions.

\section*{Standalone Adverbs}

\par\quad とても嬉しいです。\\
\par\quad I'm very happy.\\

Words like とても ("very") function as standalone, unchanging adverbs.

\section*{Adverbs Derived from Adjectives}

\par\quad 静かに話します。\\
\par\quad I speak quietly.\\

The な-adjective 静か becomes an adverb via に — the same logic we already saw with に as the anchor particle, here applied to the shift from a property into a manner.

\section*{Nouns with an Adverbial Function}

\par\quad 食べる時、手を洗います。\\
\par\quad When I eat, I wash my hands. (literally: "at the time of eating")\\

時 ("time") functions here like an adverb, but is actually a noun whose own meaning has faded considerably — its main job is to make the preceding clause attachable to a particle, much like の does in its own chapter. Words like 後 ("afterward"), 間 ("while"), 場合 ("in the case of") work on the same principle: richer variants of the same basic mechanism, each carrying its own temporal or logical meaning, but relying structurally on the same trick as の — a noun that lets a clause attach to a particle.


\chapter{あげる / くれる / もらう — Whose Eyes See the Action}

\section*{The Image}

With giving and receiving, Japanese doesn't primarily choose the verb by rank or politeness — it chooses first by whose perspective the action is told from, and whether the speaker can even take up that perspective.

\section*{The Three Basic Verbs}

\textbf{あげる — from the giver's or a neutral perspective}
\par\quad 私は友達に本をあげました。\\
\par\quad I gave my friend a book.\\

あげる tells the story from the giver's point of view — or stays neutral, without strongly stepping into either side's perspective.

\textbf{くれる — only from the receiver's perspective}
\par\quad 友達が私に本をくれました。\\
\par\quad My friend gave me a book.\\

くれる requires the speaker to take up the receiver's perspective — they have to be able to put themselves in that position. That's why くれる is used almost exclusively when the receiver is the speaker themselves (or someone very close to them): it's naturally easiest to put yourself in your own position.

\textbf{もらう — receiver's perspective, regardless of direction}
\par\quad 私は友達に本をもらいました。\\
\par\quad I received a book from my friend.\\

もらう always looks at the event from the receiver's side, no matter who's doing the giving.

\section*{Why くれる Feels So Tied to "One's Own Circle"}

Because くれる strictly requires the ability to step into the receiver's shoes, in practice it's used almost only when the receiver is the speaker or someone close to them — this isn't a fixed rule about group membership, but a natural consequence of the perspective requirement: people find it easiest to identify with their own position, or that of people close to them.

あげる, by contrast, is more flexible and can be used even when neither the giver nor the receiver is the speaker — as long as the narration stays neutral.

\section*{A Second Axis: Rank}

Within this basic direction set by perspective, there's a second layer: the specific verb changes depending on the rank relationship between giver and receiver.

\par\quad 先生に本を差し上げました。\\
\par\quad I (respectfully) gave the teacher a book.\\

\par\quad 猫にえさをやりました。\\
\par\quad I gave the cat food. (neutral to slightly condescending, common with animals, plants)\\

差し上げる is the respectful variant of あげる for higher-status receivers, やる a neutral-to-slightly-condescending variant used with animals or plants. The same logic applies in mirror image to くれる (下さる as the respectful variant) and もらう (いただく as the humble variant). Rank, then, doesn't determine the basic direction itself — only which specific word within that direction is appropriate.

\end{document}
